%!TEX root = ../main.tex
\section{Conclusion}
\label{sec:conclusion}

Quality benchmarks for scientific figure understanding tell an incomplete story. By evaluating eight vision-language models across both description quality and behavioral reliability, \textsc{SciFig-Eval} reveals that these two dimensions are largely independent. The A-R-I framework decomposes model behavior into admittance, resistance, and inductance, exposing failure modes that quality scores alone cannot detect. GPT-5.2 writes the best descriptions yet fabricates answers for blurred elements 98\% of the time, while Gemini, scoring only 1.4 MQM points lower, acknowledges visual limitations in 90\% of those same cases.

These findings carry direct consequences for deploying VLMs in scientific workflows. Selecting a model for automated figure analysis based on quality rankings alone risks embedding silent fabrication into research pipelines. The psychology-informed probes in \textsc{SciFig-Eval} offer a principled toolkit for stress-testing model honesty before deployment, grounded in decades of cognitive science on presupposition embedding, anchoring bias, and cooperative communication. The A-R-I framework itself is not specific to scientific figures and can be applied wherever models must reason under uncertainty, refuse misleading premises, or infer from partial evidence.

Several directions extend this work naturally. Broadening coverage to tables, multi-panel figures, and scientific diagrams would test whether the behavioral patterns we observe generalize across visual genres. Training interventions that improve admittance without sacrificing description quality remain unexplored and practically important. Scaling \textsc{SciFig-Eval} to additional languages would reveal whether the passive-active admittance gap persists across linguistic and cultural contexts. Finally, prompt-based strategies for closing that gap, encouraging models to express uncertainty when directly questioned, offer a low-cost path toward safer scientific deployment.

\section*{Limitations}

Our evaluation focuses on English-language bar charts, line plots, and pie charts. This scope was a deliberate design decision that enabled depth of analysis (over 11,000 evaluation instances from 250 figures), but it means our findings may not transfer to other chart types such as scatter plots, heatmaps, or schematic diagrams, nor to figures drawn from domains outside the computer science and machine learning literature prevalent on arXiv.

The 250-figure dataset, while yielding more than 23,000 model-output instances across all conditions, remains modest in absolute count. We validated stability through split-half reliability ($\rho = 0.979$) and scale analysis showing convergence at 100 figures, yet generalization to broader scientific corpora warrants further investigation. The admittance and inductance conditions were evaluated on 50 and 48 figures respectively, smaller than the other experimental conditions. Expanding these subsets is a priority for future iterations.

All automated evaluation used GPT-4o as the sole judge model. Human validation of MQM rankings yielded perfect agreement ($\rho = 1.0$), and the probe designer ablation demonstrates that behavioral scores are robust to switching the probe generator from GPT-4o to Mistral Large~3. Nevertheless, cross-judge validation with alternative evaluator models would further strengthen confidence in the scoring framework. Capability question ground truth annotations are still being finalized for complete model coverage, and preliminary results should be interpreted accordingly.

\section*{Ethics Statement}

All scientific figures in \textsc{SciFig-Eval} are sourced from arXiv preprints, which are openly accessible and licensed for research use. The dataset contains no personally identifiable information. Human annotations were performed by consenting graduate researchers who were informed of the study's purpose. The benchmark is designed to reveal behavioral limitations of vision-language models under controlled conditions, not to develop adversarial attacks or to game model performance.
